\documentclass[11pt,a4paper]{article}
\usepackage[margin=25mm]{geometry}
\usepackage{amsmath,amssymb,amsthm}
\newtheorem{lemma}{Lemma}
\newtheorem{theorem}[lemma]{Theorem}
\newcommand{\gen}[1]{\langle#1\rangle}
\title{Enumerating extensions inside centralizers}
\author{Human Compression project: finite enumeration}
\date{Working proof; distinct review required}
\begin{document}\maketitle
Retain the complete proof of enumeration by stabilizers of initial
strings, including its end intersection criterion, depth bound,
completion test and interpretation of reversal. Suppose $G\lhd A$ are
specified finite permutation groups and conjugation induces a
surjection $A\to\operatorname{Aut}(G)$. The exact identity
$|A:C_A(G)|=|\operatorname{Aut}(G)|$ verifies this last hypothesis.
The kernel of the action need not be trivial. Write $x^a=a^{-1}xa$.

For an accepted string $t=(t_1,\ldots,t_m)$, put
\[
 H=\gen t,\qquad S=C_A(t_1,\ldots,t_m),\qquad
 Q=C_A(t_1,\ldots,t_{m-1}).
\]
For the empty tuple take $S=Q=A$. The eligible next marks are precisely
the involutions in $G\cap Q$ outside $H$. This description avoids
listing every involution of $G$ before beginning the search.

\begin{lemma}\label{orbits}
Choose one representative $y$ of every $Q$-conjugacy class of
involutions contained in $G$. For each $y$, choose a complete set of
double coset representatives for $C_Q(y)\backslash Q/S$.
The elements $y^d$, over these choices of $y,d$, represent exactly
the $S$-orbits of involutions in $G\cap Q$. Discarding those in $H$
leaves exactly the eligible candidate orbits.
\end{lemma}
\begin{proof}
Normality of $G$ in $A$ makes membership in $G$ constant on every
$Q$-conjugacy class. Thus the chosen classes cover the required
involutions. For one such class, $y^{d'}=(y^d)^s$ with $s\in S$
holds if and only if $d'\in C_Q(y)dS$. Indeed the displayed equality
is equivalent to $d'(ds)^{-1}\in C_Q(y)$. This proves the asserted
bijection between double cosets and $S$-orbits. Distinct
$Q$-classes cannot merge under $S\leq Q$. Finally $S$ fixes every
mark of $t$, so it fixes $H$ pointwise. Membership in $H$ is
therefore constant on each candidate orbit, and the discard is exact.
\end{proof}

\begin{lemma}\label{closure}
Let $Y$ be the chosen class representatives in Lemma~\ref{orbits}.
The subgroup generated by all involutions of $G\cap Q$ is
\[
                         K=\gen{Y}^{\,Q},
\]
the normal closure of $\gen Y$ in $Q$. If $\gen{H,K}\ne G$,
the whole branch has no generating extension.
\end{lemma}
\begin{proof}
The involutions in question are exactly the union of the $Q$-classes
of the elements of $Y$. Their generated subgroup is therefore the
stated normal closure, and is contained in $G$. Every later mark in
a string extension commutes with $t_1,\ldots,t_{m-1}$, hence is
one of these involutions. The completed group is contained in
$\gen{H,K}$, which proves the test. Generating only $\gen Y$
would generally be insufficient; the full normal closure is required.
\end{proof}

\begin{theorem}
Replacing candidate-orbit construction in the retained enumeration
procedure by Lemma~\ref{orbits}, and its completion-subgroup test by
Lemma~\ref{closure}, records exactly one representative of each
$\operatorname{Aut}(G)$-orbit of generating string C-markings.
All ranks through $\lfloor\log_2|G|\rfloor$ suffice.
\end{theorem}
\begin{proof}
The action of $S$ is the full pointwise automorphism stabilizer of
the prefix: it is the inverse image of that stabilizer under the
surjection from $A$. Lemma~\ref{orbits} therefore supplies exactly
the candidates required by the retained enumeration theorem.
The same exact end-intersection test accepts precisely the valid
extensions. After appending $x$, the new stabilizer is $C_S(x)$,
and the new commuting group is the old $S$. These identities maintain
the invariant at every recursive call. Lemma~\ref{closure} removes
only branches that cannot generate $G$. The retained induction on
prefix length proves completeness and uniqueness, and the subgroup
index argument supplies the same depth bound. Actual representatives
may differ from the least-involution choices in the earlier program;
the represented automorphism orbits are the same.
\end{proof}

\paragraph{Implementation and attribution.}
The program computes complete conjugacy classes in $Q$, retains
the involution classes whose representatives lie in $G$, and uses
GAP's \texttt{DoubleCosetRepsAndSizes}. This operation returns a
duplicate-free complete double coset list; the sum of the returned
sizes is also checked against $|Q|$. That size check is a consistency
check and does not replace correctness of GAP's exact algorithm.
The pool for a child depends only on the parent's stabilizer and
is computed once for all that parent's children. This sharing changes
neither the candidate classes nor their normal closure. Every
generating tuple is retained literally, including disconnected
strings, and the previous exact duality test is unchanged.

The general stabilizer enumeration is the method described by
Leemans--Mixer (2012), Section~5, there attributed to
Hartley--Hulpke (2010). The double coset description is the standard
orbit--stabilizer correspondence, and normal closure is a standard
group construction. This note makes their use in this implementation
explicit; it makes no priority claim. The full preceding proof and
its cited end criterion remain required inputs. GAP group,
conjugacy, centralizer, automorphism and double coset algorithms
remain software premises. No Lean certificate or completed large
census is asserted here.
\end{document}
